\documentclass{article}

% NOTE: the official neurips_2026.sty is NOT included in this repository.
% Download it from the verified official source, media.neurips.cc/Conferences/NeurIPS2026/Styles
% (or the corresponding Overleaf/author-kit mirror), and place neurips_2026.sty
% in this same directory before compiling. Use the `preprint` option for
% anonymous / non-camera-ready submission, `final` for camera-ready.
\usepackage[preprint]{neurips_2026}

\usepackage[utf8]{inputenc}
\usepackage[T1]{fontenc}
\usepackage{hyperref}
\usepackage{url}
\usepackage{booktabs}
\usepackage{amsfonts}
\usepackage{xcolor}

\title{The Retrieval Noise Taxonomy: Disentangling Stochastic, Systematic, and Selection Effects on Self-Improvement Drift in Test-Time Continual Agents}

\author{
  Luxin Zhang \\
  \textit{[Affiliation --- to be added]} \\
}

\begin{document}

\maketitle

\begin{abstract}
Agents that improve themselves at test time by writing their own experience into memory and retrieving it on later steps sometimes get better and sometimes drift: recent work has shown that self-generated feedback alone can induce recursive degradation even when the same feedback mechanism improves agents under external supervision. We argue that part of this inconsistency is not algorithmic but architectural. The retrieval infrastructure that stores and serves an agent's memory is usually treated as a neutral component, but it introduces noise of at least two mechanistically distinct kinds: \emph{stochastic} noise from approximate nearest-neighbor search, and \emph{systematic} noise from index staleness, where a write is not yet visible to the queries that follow it. We give both a formal, checkable definition, add a third, secondary channel (\emph{selection} noise from memory eviction under a fixed budget), and hypothesize that stochastic and systematic noise can pull self-feedback drift in opposite directions: the former behaving like a regularizer, the latter compounding error by letting an agent confidently reinforce a fact it has already superseded. We describe a matched-replay ablation that holds the starting task instance fixed across every infrastructure condition, isolating infrastructure as the only source of behavioral divergence, and specify the analysis, sample size, and falsification criteria before any run is executed. We release \texttt{driftbench}, a diagnostic harness implementing all three noise channels over FAISS, tested independently of the agent loop it will eventually drive. We present this as a position paper: the contribution is the taxonomy, the causal method with its pre-registered analysis plan, and the released harness, not a completed empirical result. Real testbed integration is in progress, and we take the position that specifying what would confirm or falsify the central claim in advance is itself a contribution the field's self-improvement literature has so far lacked, independent of which outcome the eventual experiment produces.
\end{abstract}

\section{Introduction}

Consider an agent operating inside a long-horizon environment that learns, through exploration, that a particular merchant only trades ore for silver. It writes this as a memory note and moves on. Later the merchant's trade rule changes, and the agent writes a corrected note, but the correction sits in a write buffer that has not yet been merged into the searchable index, so the next query still surfaces the old rule. The agent acts on stale information, and because it is also generating its own training signal from its own actions, it writes a new note consistent with the outdated rule, reinforcing the error rather than correcting it. Now picture the same fact retrieved not through a stale index but through an approximate one: occasionally the wrong neighbor comes back, but the error is not directionally biased, and it does not compound in the same way. This scene has not been run (no experiments exist yet for this paper), but the two failure modes it stages side by side are the paper's argument in miniature: retrieval noise is not one thing.

Self-improving agents that write their own experience into memory and act on it later are the operational core of test-time continual learning \citep{skillearnbench2026}. A recent large-scale study of this setting found that self-generated feedback alone can cause recursive drift across twenty tasks and fifteen domains, while the same architecture improves reliably under external feedback \citep{skillearnbench2026}. That finding is now foundational to the field, but it leaves the mechanism underspecified: it does not say why self-feedback sometimes degrades an agent, only that it can. A separate, more applied line of work on self-improving agent architectures has begun to notice the gap directly: a recent audited skill-graph framework for agentic self-improvement lists ``infrastructure-level factors (caching, retrieval freshness, memory management)'' as explicitly unaddressed future work \citep{auditedskillgraph2026}. Nobody has yet taken up that question as a causal claim.

We do so here, and we narrow it further than ``infrastructure matters.'' Nobody has shown that retrieval noise is not one failure mode but at least two mechanistically distinct ones (stochastic and systematic) that can pull self-feedback drift in opposite directions. This is the gap the rest of the paper closes.

The paper is organized to make the taxonomy, not the ablation result, the first thing a reader encounters. Section~\ref{sec:related} surveys three adjacent literatures (the drift phenomenon itself, memory systems treated as a correctness problem, and two terminologically adjacent uses of ``collapse'' at the wrong scale for this setting) and defines the taxonomy formally against the gaps each one leaves. Section~\ref{sec:method} describes the matched-replay ablation. Section~\ref{sec:results} specifies, in advance of any run, the figure that will report the result and the criteria under which the central claim would be considered wrong. Section~\ref{sec:discussion} discusses the practical reframe the taxonomy offers regardless of which experimental outcome obtains, and Section~\ref{sec:conclusion} concludes.

We make three contributions, and we make them as a position paper rather than as a completed empirical study: the argument for the taxonomy and the method does not depend on which side of the falsification criteria in Section~\ref{sec:results} the eventual experiment lands on.

\begin{itemize}
  \item \textbf{A named taxonomy.} Stochastic, systematic, and selection noise, with formal, checkable definitions for each, intended for reuse independent of this paper's specific empirical claim.
  \item \textbf{A causal method.} A matched-replay ablation, crossing the two headline channels in a factorial design, that isolates infrastructure from task difficulty by construction rather than by statistical adjustment after the fact, together with an analysis plan and falsification criteria specified before any run is executed.
  \item \textbf{A released harness.} \texttt{driftbench} implements all three channels over FAISS and is released as its own artifact, so another group can point it at their own agent-memory stack rather than reimplementing the instrumentation from scratch.
\end{itemize}

\section{Related Work and the Retrieval Noise Taxonomy}
\label{sec:related}

The stability-plasticity tradeoff this paper's setting inherits is not new; it was established in parameter space by Kirkpatrick et al.\ \citep{kirkpatrick2017ewc} and by the continual test-time adaptation line beginning with Wang et al.\ \citep{wang2021tent} and continued by Wang et al.\ \citep{wang2022cotta}. The three literatures surveyed below inherit the same tradeoff in a retrieval-space form instead of a parameter-space one, but none of them make that translation explicit, which is the opening this section's taxonomy is built to close.

\paragraph{The drift phenomenon.} SkillLearnBench \citep{skillearnbench2026} is the empirical anchor for this paper's premise: self-feedback alone induces recursive drift across a broad task suite, while external feedback does not, and no continual-learning method wins uniformly across tasks or model scales. This result establishes that the phenomenon is real and reproducible, but the benchmark's design holds retrieval infrastructure fixed and unexamined throughout: every condition in that study runs on whatever memory backend the underlying agent framework happened to use, never as an independent variable. The paper we are building toward asks the question SkillLearnBench's design cannot answer on its own: does the infrastructure serving the memory contribute to the drift it observes.

\paragraph{Memory as a correctness problem.} A second thread treats agent memory as a systems engineering target, but frames the target as correctness rather than learning dynamics. Omri et al.\ \citep{omri2026agentmemory} give the first systems-level characterization of stateful long-horizon agent-memory workloads and formally define staleness as the number of prior sessions not yet persisted at query time, showing that several deployed memory systems accumulate staleness under realistic asynchronous consolidation schedules. This definition is adopted directly in Section~\ref{sec:method}; the paper stops at latency and freshness as systems metrics and does not connect staleness to downstream behavioral correctness or drift. Yang \citep{yang2026controlplane} studies a closely related question: how the placement of an LLM-based control plane inside a memory pipeline shapes which forgetting failure modes a system handles, evaluated across thirteen system configurations with a purpose-built adversarial benchmark. This is the closest work to the present paper found in any literature search conducted for it, and the two differ in a specific, checkable way: memory in that study is externally populated and edited through an explicit adapter protocol, with no self-generated feedback loop and no treatment of approximate search or index staleness as variables at all. The present paper's setting is active rather than passive: the agent both writes and later retrieves its own generated content, which is precisely the condition under which stale or approximate retrieval can be reinforced rather than merely encountered. Zhou et al.\ \citep{zhou2026externalization} frame memory, alongside skills and protocols, as one of several capabilities modern agents externalize from model weights into surrounding infrastructure; neither correctness-focused paper above adopts that framing explicitly, though it motivates treating memory infrastructure as load-bearing rather than incidental, which is the position this paper argues for directly.

\paragraph{Terminologically adjacent, mechanistically distant.} Two further literatures share vocabulary with this paper without sharing its mechanism or its scale. Yu et al.\ \citep{yu2026retrievalcollapse} coin ``retrieval collapse'' for the degradation of web-scale search and retrieval-augmented generation systems as AI-generated content pollutes the indexed corpus, a phenomenon operating across the open web and across many independent publishers, not inside a single agent's private memory loop. Shumailov et al.\ \citep{shumailov2024nature} show, and Dohmatob et al.\ \citep{dohmatob2025strongcollapse} formalize further, that models trained recursively on their own generated outputs accumulate distributional error across generations because each generation's synthetic data is a finite sample from the previous model's output distribution; even a small constant fraction of synthetic data in the training mix can prevent scaling from helping at all. Borji \citep{borji2024note} has since qualified the original collapse dynamics under certain data-mixing and replay conditions, though the core compounding-error mechanism motivating the analogy below is not contested. We use this literature as a motivating analogy rather than a source of a formal derivation: its central mechanism is about what \emph{fraction} of each retraining generation's signal is synthetic, not about a bias-variance distinction between noise \emph{types}, and forcing that equivalence onto approximate search versus staleness would overstate what either literature actually proves. What we borrow is narrower and, we think, defensible: an agent that retrieves and then reinforces its own self-generated memory is running a structurally similar recursive loop, one step removed from training, and it is a fair question whether an analogous asymmetry between error types shows up there too.

Across these three threads, retrieval-infrastructure noise is either assumed away as a background condition (the drift-phenomenon literature), collapsed into a single correctness problem (the memory-systems literature), or examined at a scale and mechanism that does not transfer to a single agent's private feedback loop (the corpus- and generation-level collapse literature). None decompose retrieval noise into mechanistically distinct channels or test their causal, potentially opposite-signed, effect on self-feedback drift. We propose three such channels, two of them operationalized against existing instrumentation and one proposed from scratch.

\paragraph{Stochastic noise} is the error introduced by approximate nearest-neighbor search relative to exact search over the same vectors. We operationalize it as recall@k of the approximate index against a flat index built from identical vectors, a standard, continuous, directly comparable ANN metric.

\paragraph{Systematic noise} is staleness: the number of writes that have occurred but are not yet visible to the index at the moment of a given query. We adopt this definition unchanged from Omri et al.\ \citep{omri2026agentmemory}, applying their systems-level metric as a checkable per-query quantity in an active, self-feedback setting rather than a passive logging one.

\paragraph{Selection noise}, kept secondary throughout this paper, is the information loss from eviction under a fixed memory budget. Unlike the other two channels, its metric has no precedent in the literatures surveyed above: we propose measuring it as the retention rate of a budgeted memory against an unbounded-memory oracle run of the same task instance, checking at each query point whether the budgeted agent's memory still contains what an unconstrained counterpart would. We flag this metric as novel rather than adopted, and for that reason treat selection as an appendix-tier ablation rather than a headline channel on equal footing with the other two.

\section{Method}
\label{sec:method}

\subsection{Matched-replay ablation}

The central methodological risk in attributing drift to infrastructure rather than to task difficulty is confounding: harder, longer-horizon tasks plausibly force smaller retrieval budgets, staler indices, and more aggressive eviction simultaneously, so an observed correlation between ``worse infrastructure'' and ``more drift'' could reflect nothing more than ``harder tasks'' acting on both. We rule this out by construction rather than by statistical adjustment. Every grid condition replays from the \emph{same starting task instance}, the same seed and the same initial world state, so infrastructure is the only variable that can differ between conditions at the moment of divergence. Trajectories are not held fixed after that point; once retrieval begins returning different results under different conditions, the agent's own actions diverge, which is the entire mechanism under study rather than a confound in it.

\subsection{Experimental grid}

The ablation is a two-way factorial crossing the stochastic and systematic channels, three levels each, nine cells, replicated on two testbeds. A full three-way factorial including selection was considered and rejected: it multiplies cost without answering a question the paper needs to answer, since selection's role is secondary and its metric is untested against precedent. Instead, selection is varied one-at-a-time against the shared exact-search, fresh-index baseline cell, which doubles as the grid's control condition. Concretely, the stochastic channel is swept across exact (flat) search, moderate approximation (HNSW with \texttt{ef\_search=16}), and aggressive approximation (HNSW with \texttt{ef\_search=4}); the systematic channel is swept across an index rebuilt every step, a moderate rebuild cadence, and a stale cadence in which many writes accumulate before a rebuild. The full grid specification is versioned in the released harness (\texttt{configs/grid\_default.yaml}).

The factorial design over the two headline channels exists specifically to answer a question one-at-a-time sweeps cannot: does the effect of staleness depend on how approximate search already is, or vice versa. This is the same interaction term that would need to be reported to defend an ``opposite-signed'' claim against the confounding objection above, so the grid's shape is not incidental to the analysis but required by it.

\subsection{Analysis}

Because every cell in the factorial grid replays the same starting instance, the comparison across cells for a given instance is paired rather than independent. We analyze the design with a repeated-measures / mixed-effects model, treating task instance as a random effect and the two channel levels, together with their interaction, as fixed effects, fit separately per testbed rather than assuming the effect transfers between them. The primary reported quantity is the stochastic-by-systematic interaction term, with a confidence interval rather than a bare significance threshold: at the sample size available to a workshop-scale study, threshold-based framing would overstate the precision the design can actually deliver, and an effect-size estimate with an honest interval is the more defensible claim to make from it.

We fix $N = 15$--$20$ task instances per cell. No pilot data exists to justify this number through a power calculation, and we do not pretend otherwise; the number reflects a compute budget rather than a target power level, and results are reported and interpreted accordingly. Across the nine-cell factorial grid and two testbeds this yields approximately 324 matched-replay rollouts, plus roughly 70--100 additional rollouts for the selection channel's one-at-a-time conditions, for a total near 400.

\subsection{Base model, embedding, and inference stack}

Every rollout uses Llama-3.1-8B-Instruct as the base agent model, served through vLLM, with BAAI/bge-base-en-v1.5 as the embedding model for all memory vectors. This choice is not arbitrary: AgentOdyssey's own evaluation already validates both Llama-3.1-8B and the smaller Qwen3-4B as open-weight agents in this exact class of environment, so neither model nor stack is introduced without precedent \citep{zhang2026agentodyssey}. Evo-Memory's own evaluation, by contrast, used only closed-weight API models \citep{wei2025evomemory}; using an open-weight model there means results on that testbed will not be directly comparable to Evo-Memory's published baselines, which we treat as an acknowledged limitation rather than a blocker, since that testbed is used here as a second matched-replay stage for this paper's own ablation rather than as a leaderboard to be beaten. AgentOdyssey additionally reports that long-context agents incur quadratic token cost over long horizons while fixed-size-memory and retrieval-augmented agents remain linear \citep{zhang2026agentodyssey}; the fixed-size-memory design used throughout this paper is required by that finding as much as it is a methodological preference, since a quadratic-cost design would make roughly 400 long-horizon rollouts intractable within a workshop compute budget regardless of infrastructure condition.

Every episode is bounded by a fixed step count rather than by wall-clock time. This is a design requirement, not an incidental detail: aggressive approximate search is faster than exact search, and under a wall-clock budget the more-approximate conditions could complete more steps than the more-exact ones, confounding retrieval noise with effective exploration budget rather than isolating it.

The generation model itself samples stochastically (temperature 0.7), which is a second source of run-to-run variance distinct from the infrastructure noise under study, and matched-replay does not by itself separate the two. We do not treat this as a confound requiring a separate control, for a specific reason: the mixed-effects analysis in Section~\ref{sec:method} already models task instance as a random effect, and sampling variance is nested inside that same instance-level term regardless of which grid cell produced it, so it inflates the residual rather than biasing the interaction estimate the paper reports. This is weaker than an explicit multi-sample control per cell would be, and we say so rather than assume it away: a reviewer could reasonably ask for multiple LLM samples per condition to isolate sampling variance from infrastructure variance directly, and we treat that as a natural extension of the design rather than as a threat that has already been resolved.

\subsection{Evaluation protocol}

We adopt the recall-versus-proxy protocol from Dong et al.\ \citep{dong2026seqmemeval} rather than designing an evaluation scheme from scratch. At each checkpoint, we record both a proxy metric (downstream task success) and a free-form recall score obtained by removing supporting context and asking the agent to answer from memory alone. The gap between these two numbers is itself diagnostic: a companion critique of test-time training methods found that loss-based proxies can improve while free-form recall remains at zero, meaning a memory method can appear to be working while retaining nothing checkable \citep{song2026beyondperplexity}. Reporting both numbers at every checkpoint, rather than only the proxy metric, is what allows this failure mode to be visible in the results if it occurs. We additionally track a drift metric adapted from Mishra \citep{mishra2026rdumbplusplus}, computing the KL divergence of each episode's action or output distribution against a fixed early-episode reference, as a model-agnostic complement to the recall-versus-proxy gap.

\subsection{Harness}

All of the above is implemented in \texttt{driftbench}, released alongside this paper at \url{github.com/luxinlabs/Neurips_ttcl2026}. The harness composes the three noise channels (an approximate-search wrapper built on FAISS, a shadow-index writer that defers pending writes until a configurable rebuild cadence, and a budgeted-memory wrapper supporting FIFO, LRU, and importance-weighted eviction) into a single memory store exposing all three checkable quantities from a single retrieval call. A generic self-feedback episode loop (observe, retrieve, act, reflect, write) is implemented against injectable model, embedding, and testbed interfaces, so the loop itself is fully tested with deterministic stand-ins independent of which base model or testbed backs it. As of this submission, the core instrumentation (the three noise-channel wrappers, the grid builder, the matched-replay orchestrator, the recall-versus-proxy checkpointing, and the drift metric) is implemented and covered by forty passing unit tests. The two testbed adapters (AgentOdyssey, Evo-Memory) are explicit, documented stubs conforming to the same interface; real integration against each benchmark's public repository is in progress and is the only remaining component standing between this harness and the results reported in Section~\ref{sec:results}.

\section{Results}
\label{sec:results}

As a position paper, this section's contribution is the specification itself: what the results will report, how they will be visualized, and how they will be judged, fixed in advance of any run so the analysis cannot be shaped after the fact by whatever pattern the data happens to show. No experimental results exist at the time of this submission; real testbed integration against AgentOdyssey and Evo-Memory is in progress, described in Section~\ref{sec:method}.

\paragraph{Headline figure.} The primary result will be reported as a two-panel figure. Panel A is a marginal-effects plot: two lines with confidence-interval ribbons, showing the drift metric against each channel's level independently, holding the other channel at its baseline. This is the most direct visual rendering of an opposite-signed claim, and is designed to be legible without reading a table. Panel B is the full $3\times3$ heatmap across both channels, with the interaction term's estimate and interval annotated directly on it, carrying the robustness argument against the task-difficulty confound described in Section~\ref{sec:method}. Both testbeds will be shown as side-by-side small multiples rather than pooled, since the mixed-effects analysis is fit separately per testbed and the figure should not pre-commit to pooling before the analysis does.

\paragraph{Falsification criteria, stated in advance.} A null interaction term (a confidence interval that includes zero) is not evidence against the paper's central claim. It would mean the two channels act additively and independently of each other, which is the best-case robustness result: it answers the task-difficulty confound directly without touching the opposite-signed claim at all. Two other outcomes would count as genuine evidence against the paper as framed. If both channels show a detectable main effect but in the \emph{same} direction, the specific ``opposite-signed'' claim is falsified; each channel could still be real and mechanistically distinct, but the taxonomy's citable hook is gone, and the paper would need to be reframed as a different, less distinctive contribution. If \emph{neither} channel shows a detectable main effect at all, the foundational premise borrowed from SkillLearnBench does not reproduce under this design, which would call the whole approach into question rather than any single contribution within it.

\paragraph{The one-channel fallback.} Suppose only one channel shows a detectable effect while the other does not, most plausibly the systematic one, given its structural resemblance to the compounding-error mechanism in the model-collapse literature discussed in Section~\ref{sec:related}. In that case the paper's first contribution would be reframed from ``two opposite-signed channels'' to ``the taxonomy correctly predicts which noise type is behaviorally consequential and which is not.'' This is arguably a more directly useful claim for a practitioner than the stronger opposite-signed version would have been: it tells them specifically where to spend engineering effort, index freshness rather than search accuracy, for instance, rather than leaving ``retrieval quality'' as a single undifferentiated concern. The paper's title does not commit to the opposite-signed outcome, and the model-collapse analogy in Section~\ref{sec:related} would if anything be strengthened by this fallback, since the mechanism it describes is specifically a systematic-noise phenomenon.

\section{Discussion}
\label{sec:discussion}

If the strong form of the central claim holds, the practical implication for anyone deploying a self-improving agent is that ``retrieval quality'' is not a single dial to be turned toward ``more accurate.'' An engineering decision to use exact rather than approximate search and a separate decision about how often to rebuild an index are answers to two different risks, and tuning them as though they were one setting means at best wasting effort and at worst tuning in the wrong direction for one of the two. If instead only the systematic channel proves consequential, the practical guidance narrows but does not disappear: it says specifically that staleness, not search approximation, is the risk worth engineering around in a self-feedback setting, which is still a more actionable claim than either ``retrieval quality matters'' or no claim at all.

The paper's scope leaves several things unaddressed deliberately rather than by oversight. The taxonomy's third channel, selection, is instrumented and tested but carries a metric with no precedent in prior work, and is treated throughout as secondary evidence rather than a claim on equal footing with the other two. The model-collapse analogy in Section~\ref{sec:related} is used to motivate a hypothesis, not to derive one; a formal treatment connecting compounding synthetic-data error to compounding retrieval error remains future work, and we think the results reported here, particularly under the one-channel fallback, would motivate that treatment more precisely than attempting it without evidence first.

\section{Conclusion}
\label{sec:conclusion}

Retrieval infrastructure has generally been treated as a neutral storage layer sitting underneath an agent's memory, rather than as a source of noise with its own behavioral consequences. This paper argues it is at least two mechanistically distinct noise channels, with different, possibly opposite, consequences for whether a self-improving agent's behavior drifts under its own feedback. Whichever of the outcomes specified in Section~\ref{sec:results} obtains, the payoff is meant to be decision-relevant rather than merely descriptive: a practitioner should come away knowing which infrastructure knob is worth their engineering effort, not only that infrastructure matters in some unspecified way. The taxonomy, the matched-replay causal method, and the released harness are offered as reusable infrastructure for a question the field has so far only named as future work.

Beyond the immediate empirical result, three directions follow naturally. The taxonomy could be extended to noise channels neither testbed used here captures, distributed or networked consolidation delay under multi-replica deployment being the most obvious. A diagnostic that measures drift under a given infrastructure configuration invites a closed-loop counterpart: an adaptive policy that detects drift and adjusts staleness or approximation tolerance in response, rather than fixing them in advance. And the model-collapse analogy this paper deliberately declined to formalize is, we think, worth formalizing properly once real evidence exists to formalize it around.

\section*{Limitations}

The selection channel's retention-rate metric is proposed in this paper rather than adopted from prior work, and is flagged as such rather than presented with the same evidentiary weight as the stochastic and systematic channels' metrics, both of which are adopted from existing instrumentation. The Evo-Memory testbed's own evaluation used only closed-weight models, so results obtained there with an open-weight base model are not directly comparable to its published baselines. The sample size, fifteen to twenty task instances per cell, was set by compute budget rather than by a power calculation grounded in pilot data, and results are reported with confidence intervals rather than significance thresholds for that reason. The model-collapse analogy motivating part of the paper's framing is explicitly not a formal derivation, and a published critique of the originating result exists and is discussed in Section~\ref{sec:related} rather than omitted. Every rollout in this study uses a single base model, Llama-3.1-8B-Instruct; whether the reported effects generalize across model scale is untested. The design also does not separate sampling variance in the generation model from infrastructure-induced variance through an explicit multi-sample control per cell, relying instead on task instance as a shared random effect covering both, discussed in Section~\ref{sec:method}. Finally, and most significantly, the paper's central empirical claim has not yet been tested: the harness described in Section~\ref{sec:method} is implemented and unit-tested, but real integration against the AgentOdyssey and Evo-Memory testbeds remains in progress, and the results in Section~\ref{sec:results} are a specification rather than a finding.

\section*{Data Availability Statement}

All code implementing the taxonomy's three noise channels, the matched-replay orchestrator, the evaluation protocol, and the experimental grid configuration is publicly available at \url{github.com/luxinlabs/Neurips_ttcl2026} under the harness name \texttt{driftbench}. Testbed adapters for AgentOdyssey and Evo-Memory are present in the same repository as documented stubs pending integration; no experimental result data exists at the time of this submission, per the Limitations section above.

\section*{Ethics Statement}

This work does not involve human subjects, personally identifiable data, or any data collection beyond procedurally generated text-game environments and existing public benchmarks. No ethical review was required.

\section*{Author Contributions}

Sole author. Luxin Zhang conceived the taxonomy, designed the matched-replay ablation and evaluation protocol, and implemented the \texttt{driftbench} harness.

\section*{Conflict of Interest Statement}

The author declares no conflicts of interest.

\section*{Funding}

This work received no external funding. Computational resources for the planned experiments are expected to draw on GPU credits provided through the NeurIPS 2026 TTCL workshop's compute sponsorship (Lambda); no funds were received in connection with the preparation of this manuscript.

\section*{Use of AI Tools}

The author used Claude Sonnet 5 (Anthropic) throughout the preparation of this manuscript. Specifically, the tool was used for literature search and citation verification (author lists, venues, and arXiv identifiers were fetched directly from source pages rather than relied upon from model memory), drafting of the manuscript text following a human-directed Socratic planning process in which the author made all substantive research and framing decisions, implementation and unit testing of the accompanying \texttt{driftbench} code harness, an internal simulated peer-review pass used to identify and address weaknesses prior to submission, and conversion of the manuscript to NeurIPS-format \LaTeX. The AI assistant did not perform any data analysis: no experimental results exist at the time of this submission, per the Limitations section above. All AI-assisted output was reviewed and verified by the author, who takes full responsibility for the content of this manuscript.

\emph{Placement note: per NeurIPS's AI-usage policy (accessed 2026-04-09, \texttt{neurips.cc/public/EthicsGuidelines}), this disclosure belongs in the Acknowledgements section or a dedicated ``Use of AI Tools'' subsection before References --- placed here as the latter. The policy snapshot is over four months old at the time of this draft; verify against the current NeurIPS page before final submission.}

\bibliographystyle{plainnat}
\bibliography{references}

\end{document}
